\chapter{Subgraph matching: el problema}
\label{cap:subgraph-matching}

\begin{intuicion}
\emph{Subgraph matching} es: dado un grafo grande \(G\) (el dato)
y un grafo chico \(P\) (el patrón o hipótesis), encontrar todas las
formas de \emph{embeber} \(P\) dentro de \(G\). En GraphHunter el
``patrón'' es la hipótesis del analista; \(G\) es el grafo temporal
ingestado. El problema en su formulación más general es NP-duro
---el famoso \emph{subgraph isomorphism problem}--- pero en la
práctica se vuelve tractable por tres razones combinadas: las
hipótesis son chicas (\(k \leq 6\) saltos), las etiquetas
(\(\mathrm{type}_V\), \(\mathrm{type}_E\)) podan agresivamente, y
las restricciones temporales acotan la \emph{ventana de búsqueda}.
Este capítulo enuncia las variantes formales (decisión, conteo,
enumeración), distingue isomorfismo de homomorfismo, y muestra por
qué el motor implementa enumeración.
\end{intuicion}

\section{Tres niveles de pregunta}

Cuando un analista escribe una hipótesis, puede querer una de tres
respuestas distintas:

\begin{enumerate}
    \item \textbf{Decisión}: ``¿Existe al menos un walk de hipótesis
    en \(G\)?'' \\
    Respuesta: \texttt{true}/\texttt{false}.
    \item \textbf{Conteo}: ``¿Cuántos walks de hipótesis hay?'' \\
    Respuesta: un entero.
    \item \textbf{Enumeración}: ``Listamelos uno por uno.'' \\
    Respuesta: una secuencia (potencialmente larga) de tuplas de
    vértices.
\end{enumerate}

GraphHunter implementa la versión 3, \textbf{enumeración}, con un
\emph{cap} configurable (default 10\,000 resultados) que la trunca
para evitar que un \emph{hub} explote la salida. ¿Por qué la
enumeración? Porque el analista no quiere saber cuántos atacantes
hay ---quiere saber \emph{cuáles} (qué usuario, qué host, qué IP).
La salida es la entrada del próximo paso de la investigación.

\begin{nota}
Las versiones de decisión y conteo son sub-problemas de enumeración
pero, sorpresa: en general son \emph{tan caras} como enumeración.
Para preguntas sin truncate, el motor tiene que generar candidatos y
verificarlos uno por uno; saber ``hay al menos uno'' suele requerir
generar al menos uno. Hay casos donde decisión es estrictamente más
fácil (cuando hay índices que la responden directo: NLF, k-hop
reachability), pero esos son optimizaciones, no soluciones generales.
Las desarrollamos en los capítulos \ref{cap:prune-nlf} y
\ref{cap:prune-khop}.
\end{nota}

\section{Patrón \(P\) versus dato \(G\)}

\begin{definition}[Patrón]
Un \emph{patrón} (en GraphHunter, una \texttt{Hypothesis} con
\emph{shape} \emph{chain}) de longitud \(k\) es una tupla
\((P, \mathrm{type}_V^P, \mathrm{type}_E^P)\) donde \(P =
(V_P, E_P)\) es un grafo dirigido etiquetado, \(|V_P| = k+1\),
\(|E_P| = k\), y \(P\) es una cadena: hay un orden total
\(p_0, p_1, \dots, p_k\) tal que \(E_P = \{(p_0, p_1), (p_1, p_2),
\dots, (p_{k-1}, p_k)\}\).

Para \emph{shape} \emph{graph}, \(P\) puede ser cualquier grafo
\(\alpha\)-acíclico (capítulo~\ref{cap:hipergrafos}), no sólo una
cadena.
\end{definition}

El \emph{dato} es el grafo temporal etiquetado \(G\) construido por
la ingesta. Las dimensiones típicas:

\begin{center}
\renewcommand{\arraystretch}{1.2}
\begin{tabular}{ll}
\toprule
\(|V_P|, |E_P|\) & \(\leq 6\) y \(\leq 6\) (hipótesis del analista) \\
\(|\V|, |\E|\) & hasta \(10^7\) y \(10^8\) (grafo del motor) \\
\bottomrule
\end{tabular}
\end{center}

Esa asimetría dramática es la que hace que el problema sea ---a pesar
de la dureza teórica--- tractable en la práctica.

\section{Embeddings: los matches del patrón}

¿Qué significa que \(P\) ``aparezca'' en \(G\)? Hay dos definiciones
clásicas, y la elección entre ellas tiene consecuencias.

\subsection{Homomorfismo}

\begin{definition}[Homomorfismo de grafos]
Un \emph{homomorfismo} de \(P = (V_P, E_P)\) en \(G = (\V, \E)\)
(ambos etiquetados, dirigidos) es una función
\(\sigma \colon V_P \to \V\) tal que:
\begin{enumerate}
    \item Para todo \(p \in V_P\), \(\mathrm{type}^P(p) =
    \mathrm{type}(\sigma(p))\).
    \item Para todo \((p, q) \in E_P\), existe una arista
    \((\sigma(p), \sigma(q)) \in \E\) con
    \(\mathrm{type}_E((\sigma(p), \sigma(q))) =
    \mathrm{type}_E^P((p, q))\).
\end{enumerate}
\end{definition}

\noindent
La función \(\sigma\) puede mapear distintos vértices del patrón al
\emph{mismo} vértice de \(G\). Es decir: el embedding puede
``colapsar''.

\subsection{Isomorfismo (subgraph isomorphism)}

\begin{definition}[Embedding isomórfico]
Un \emph{embedding isomórfico} es un homomorfismo \(\sigma\) que
además es \textbf{inyectivo}: \(p \neq q\) implica \(\sigma(p) \neq
\sigma(q)\).
\end{definition}

\subsection{¿Cuál es la convención de GraphHunter?}

GraphHunter usa una variante intermedia: la \(k_s\)-\emph{simplicidad}
del capítulo~\ref{cap:grafos-temporales}. Por defecto, \(k_s = 1\)
(ningún vértice se repite, equivale a embedding isomórfico).
Configurar \texttt{\{k=2\}} en la DSL afloja: cada vértice puede
aparecer hasta dos veces. \texttt{\{k=k+1\}} (donde \(k\) es la
longitud de la hipótesis) lo torna en homomorfismo puro.

\begin{advertencia}
Esta diferencia importa muchísimo. Considerá la hipótesis
\texttt{User -[Auth]-> Host -[Auth]-> User}. Bajo isomorfismo, las
respuestas son ternas \((u_1, h, u_2)\) con \(u_1 \neq u_2\). Bajo
homomorfismo, la terna \((u_1, h, u_1)\) ---el mismo usuario
``volviendo''--- también cuenta. Para detectar \emph{lateral
movement} suele importar la primera; para detectar \emph{credential
abuse}, la segunda.
\end{advertencia}

\section{Por qué el problema es duro en general}

\begin{theorem}[Cook-Levin via Clique]
El problema de decisión asociado a subgraph isomorphism (``¿hay un
embedding isomórfico de \(P\) en \(G\)?'') es NP-completo, incluso
si \(P\) y \(G\) son grafos no dirigidos sin etiquetas.
\end{theorem}

\begin{proof}[Idea de la prueba]
Reducción desde \emph{clique}. Dado \((G, k)\) ---instancia de
clique: ¿tiene \(G\) un clique de tamaño \(k\)?--- construimos
\(P = K_k\) (el grafo completo de \(k\) vértices) y preguntamos si
\(P\) embebe isomórficamente en \(G\). Existe embedding sii \(G\)
tiene un \(k\)-clique. \(\square\)
\end{proof}

Esto significa que ningún algoritmo conocido resuelve subgraph
matching en tiempo polinomial \emph{en general} (asumiendo \(P \neq
NP\)). En el peor caso, el motor está obligado a explorar un espacio
combinatorio del tamaño de \(\binom{n}{k}\) para hipótesis de
\(k+1\) vértices.

\subsection{Por qué entonces es tractable en la práctica}

Tres palancas que GraphHunter usa todas a la vez:

\begin{enumerate}
    \item \textbf{Etiquetas \(\mathrm{type}_V\) y \(\mathrm{type}_E\)
    podan masivamente.} El primer paso del DFS sólo considera
    vértices del tipo \(L_0^V\). Si la hipótesis es
    \texttt{User -[Auth]-> Host -[Connect]-> IP} y el grafo tiene
    \(10^6\) \texttt{User}s pero \(10^4\) \texttt{Host}s y \(10^5\)
    \texttt{IP}s, el factor \(10^6 \times 10^4 \times 10^5 = 10^{15}\)
    naive se reduce a las parejas que \emph{sí} tienen aristas con la
    etiqueta correcta ---típicamente \(10^7\) órdenes de magnitud
    menos.
    \item \textbf{La \emph{ventana temporal}} \(\Delta t\) recorta el
    sub-slice de \(N^{+}(v)\) que el matcher considera. En vez de
    iterar \(\outdeg(v)\) aristas, itera las que caen en
    \([\tau_{i-1}, \tau_{i-1} + \Delta t)\). Para \(\outdeg = 10^4\)
    aristas distribuidas en una semana y \(\Delta t = 30\) min,
    eso son \(\sim 30\) aristas por vértice intermedio.
    \item \textbf{Los \emph{prunes}} (NLF, k-hop reachability)
    descartan \emph{starts} antes de iniciar el DFS. Un vértice
    cuyos vecinos no tienen el tipo correcto del primer paso, o
    desde el cual no hay \emph{ningún} camino al \texttt{dest\_type}
    final en \(\leq \texttt{MAX\_DEPTH}\) saltos, jamás genera
    resultados. Se descarta sin tocar el DFS.
\end{enumerate}

\section{Variante: subgraph matching temporal}

La definición clásica de subgraph isomorphism no incluye el tiempo.
GraphHunter agrega la restricción del capítulo
\ref{cap:grafos-temporales}: cada embedding debe ser \emph{temporalmente
válido} y \emph{\(\Delta t\)-acotado}.

\begin{definition}[Match temporal]
Un \emph{match temporal} de un patrón \(P\) (chain de longitud \(k\))
en un grafo temporal \(G\) es un par \((\sigma, \mathbf{e})\) donde:
\begin{itemize}
    \item \(\sigma \colon V_P \to \V\) es una asignación que respeta
    etiquetas de vértice y \(k_s\)-simplicidad.
    \item \(\mathbf{e} = (e_0, e_1, \dots, e_{k-1})\) es una secuencia
    de aristas con \(e_i = (\sigma(p_i), \sigma(p_{i+1}), \tau_i)\),
    \(\mathrm{type}_E(e_i) = \mathrm{type}_E^P((p_i, p_{i+1}))\).
    \item \(\tau_0 < \tau_1 < \dots < \tau_{k-1}\) (monotonicidad).
    \item \(\tau_{k-1} - \tau_0 \leq \Delta t\) (ventana).
    \item Cualquier predicado adicional sobre \(\sigma\) o
    \texttt{metadata} se cumple.
\end{itemize}
\end{definition}

\noindent
Notar que \(\mathbf{e}\) es parte del match: el mismo \(\sigma\)
puede tener varios matches temporales si las aristas multi-grafo
(con \(\tau\) distintos) se eligen distinto.

\section{Operadores: qué algoritmo usa el motor}

GraphHunter no implementa \emph{un} algoritmo de subgraph matching:
implementa un \emph{zoológico} de operadores y un \emph{planner} que
elige uno. La Parte~V del libro los desarrolla; acá los nombramos para
que el lector tenga el mapa:

\begin{itemize}
    \item \textbf{DFS recursivo con \emph{prunes}}
    (capítulo~\ref{cap:dfs-baseline}): el \emph{baseline}; funciona
    para cualquier hipótesis. Su complejidad de peor caso es
    \(O(m \cdot \deg^k)\); los \emph{prunes} (NLF, k-hop) la bajan
    drásticamente en grafos reales.
    \item \textbf{LFTJ-K\(_4\)} (capítulo~\ref{cap:lftj}): un
    \emph{worst-case-optimal join} para hipótesis cuya \emph{shape}
    es \(K_4\) (cuatro vértices, las seis aristas dirigidas). Logra
    el AGM bound de \(O(m^2)\) en lugar de \(O(m \cdot \deg^3)\).
    \item \textbf{Yannakakis chain}
    (capítulo~\ref{cap:yannakakis-chain}): para cualquier hipótesis
    \(\alpha\)-acíclica con shape chain, una semi-join reducer en
    dos pasadas eliminando \emph{rows} sin futuro, y enumeración al
    final.
    \item \textbf{Yannakakis tree-walk}
    (capítulo~\ref{cap:yannakakis-tree}): el caso general
    \(\alpha\)-acíclico (estrellas, árboles, hipergrafos) con
    enumeración recursiva sobre el \emph{join tree}.
\end{itemize}

El \emph{planner} (capítulo~\ref{cap:planner}) decide cuál
dispatchar inspeccionando la \emph{shape} de la hipótesis. Por
defecto cae al DFS; las hipótesis con propiedades especiales
(\(K_4\) homogéneo, \(\alpha\)-aciclidad) ganan operadores más
agresivos.

\section{Por qué el AGM bound importa}

El capítulo~\ref{cap:joins-y-grafos} deriva el AGM bound desde cero,
pero adelantamos la idea: para una hipótesis \(P\), existe una
cantidad \(\rhostar(P) \in [1, |E_P|]\) tal que el número de matches
posibles es a lo sumo \(O(m^{\rhostar(P)})\). Esa cantidad se llama
el \emph{fractional edge cover number} y depende sólo de la
estructura del patrón.

\begin{itemize}
    \item Para una \emph{cadena} de \(k\) aristas, \(\rhostar = k\).
    El DFS naive ya alcanza ese bound; no hay mucho margen.
    \item Para un \emph{triángulo}, \(\rhostar = 1.5\). El DFS
    escala como \(O(m \cdot \deg) = O(m \cdot \frac{m}{n})\) en el
    peor caso ---\emph{peor} que \(O(m^{1.5})\). Hay margen
    significativo.
    \item Para \(K_4\), \(\rhostar = 2\). El DFS es \(O(m \cdot
    \deg^3)\); LFTJ alcanza \(O(m^2)\). En grafos densos eso es
    varios órdenes de magnitud.
\end{itemize}

\noindent
Esto es lo que justifica formalmente por qué el motor tiene un
operador especializado para \(K_4\) y no, digamos, para cadenas: el
\emph{lift} algorítmico está disponible donde la estructura del
patrón lo permite.

\section{Resumen}

\begin{itemize}
    \item Subgraph matching es: dado \(G\) y \(P\), encontrar
    embeddings de \(P\) en \(G\). En GraphHunter, la variante
    operativa es \textbf{enumeración con cap}, sobre \textbf{matches
    temporales} (monotónicos y \(\Delta t\)-acotados), respetando
    \textbf{etiquetas} y \textbf{\(k_s\)-simplicidad}.
    \item El problema de decisión es NP-completo en general
    (reducción desde clique).
    \item En la práctica es tractable por la combinación de
    etiquetas, ventana temporal, prunes (NLF, k-hop), y la elección
    inteligente de operador físico (DFS, LFTJ, Yannakakis).
    \item El \(\rhostar\) del patrón ---el \emph{fractional edge
    cover number}--- determina cuánto \emph{lift} sobre el DFS naive
    es teóricamente posible. Para patrones donde \(\rhostar < k\),
    los \emph{worst-case-optimal joins} ganan; para cadenas, el
    DFS bien-podado es ya cerca del óptimo.
\end{itemize}

\begin{ejercicios}
\begin{enumerate}
    \item \starred{} Para el patrón ``triángulo'' (3 vértices, 3
    aristas dirigidas formando un ciclo), enumerá todos los
    embeddings isomórficos en el grafo de la
    figura~\ref{fig:gt-ejemplo} del capítulo anterior. ¿Cuántos hay?
    ¿Y bajo homomorfismo (permitiendo colapsar vértices)?

    \item Demostrá por reducción que el problema ``¿existe un
    Hamiltonian path en \(G\)?'' se puede resolver con un solo subgraph
    matching. (Pista: ¿qué patrón \(P\) usás?)

    \item \starred{} La \emph{cadena} de longitud \(k\) tiene
    \(\rhostar = k\). Dada esta propiedad, ¿qué predice el AGM bound
    sobre el número máximo de matches en un grafo de \(m\) aristas?
    Compará con el bound trivial \(O(\deg^k)\).

    \item Considerá una hipótesis con shape ``estrella'': un vértice
    central conectado a \(k\) vértices distintos vía aristas
    salientes. ¿Cuál es \(\rhostar\) de esa estructura? (Pista:
    pensá cuál es el \emph{minimum fractional edge cover}.)

    \item ¿Por qué decidir si dos grafos arbitrarios son isomorfos
    (\emph{graph isomorphism problem}) NO es lo mismo que subgraph
    isomorphism? ¿Cuál es ``más fácil'' (en el sentido de la teoría
    de complejidad)?

    \item Imaginá una hipótesis donde GraphHunter devuelve 9.999
    resultados (justo bajo el cap de 10.000). ¿Cómo sabés si los
    9.999 son ``todos los matches'' o ``sólo los primeros 10K de
    muchos más''? Diseñá una API que distinga los dos casos.
\end{enumerate}
\end{ejercicios}
